% The geometric Hurst exponent of the crests of the French Alps and Pyrenees.
% Figures in figs/ (code/fig_paper.py), numbers in data/ (code/hgeo.py).
% Compile with `make' in this folder (pdflatex twice, for the references).
\documentclass[11pt,a4paper]{article}
\usepackage[margin=2.3cm]{geometry}
\usepackage[T1]{fontenc}
\usepackage{lmodern}
\usepackage{microtype}
\usepackage{amsmath,amssymb,amsthm}
\usepackage{graphicx}
\usepackage{booktabs}
\usepackage{array}
\usepackage{tabularx}
\usepackage{xcolor}
\usepackage[font=small,labelfont=bf]{caption}
\usepackage[round,authoryear]{natbib}
\usepackage[colorlinks=true,linkcolor=blue!50!black,citecolor=blue!50!black,
            urlcolor=blue!50!black]{hyperref}
\graphicspath{{figs/}}
\renewcommand{\topfraction}{0.9}
\renewcommand{\bottomfraction}{0.6}
\renewcommand{\textfraction}{0.1}
\renewcommand{\floatpagefraction}{0.75}
\newcommand{\Hgeo}{H_{\mathrm{geo}}}
\newcommand{\Hroute}{H_{\mathrm{route}}}
\newcommand{\Hdyn}{H_{\mathrm{dyn}}}
\newcommand{\code}[1]{\texttt{#1}}

\title{How straight is a ridge?\\
The geometric Hurst exponent of the crests of the French Alps and Pyrenees}
\author{Matteo Di Sante}
\date{Draft, 12 November 2026}

\begin{document}
\maketitle

\begin{abstract}
A paraglider that flies cross-country in the mountains follows the ridges, and
its track inherits their shape. Before the Hurst exponent of such a track can
say anything about how pilots fly, one must know how much of it the terrain
alone sets. We give that baseline. We extract the crest lines of the French
Alps and Pyrenees from the 25~m IGN elevation model, join the pieces into
long crest paths, and let a ball run along them at constant speed. The
time-averaged mean squared displacement of the ball, fitted from 100~m (four
grid cells) up, gives a geometric Hurst exponent $\Hgeo = 0.922$
(0.918--0.929) in the Alps and $0.916$ (0.910--0.931) in the Pyrenees. The
exponent is not a constant: it falls from about 0.98 at 100~m to 0.94 at
1~km and 0.87 beyond 10~km, along the same curve in both massifs. Box
counting, an independent geometric measure, agrees: $1/D$ follows the local
$H$ to within 0.01--0.03 up to about 10~km. A flight with $H$ close to these
values is therefore as straight as the ridges it follows, and no straighter.
We give the definitions, data and code needed to put real flights on the
same scale, and the decomposition (terrain, route, time) the comparison
calls for.
\end{abstract}

\section{Introduction}

Cross-country paragliding in the mountains is ridge flying. Thermals rise
from the sunny slopes and the lift is strongest close to the crests, so a
pilot who wants to go far follows the ridges, climbs in a thermal, glides
to the next one, and so on. The GPS track of a flight records two things at once:
the shape of the terrain and the way the pilot moves through it, with
straight glides, detours and long pauses spent circling in thermals.

A common way to describe such a track is the mean squared displacement and
its exponent. If $\delta^2(\tau) \sim \tau^{2H}$, a Hurst exponent $H=1$ means
a ballistic, straight motion and $H=1/2$ a diffusive one. Measured on a flight,
$H$ mixes the two ingredients. A pilot who flies perfectly straight at constant
speed along a winding ridge has $H<1$, and the shortfall is the terrain's, not
the pilot's. The question of this paper is therefore simple: \emph{how much of
the Hurst exponent of a ridge flight can the terrain alone explain?}

To answer it we need the Hurst exponent of the crests themselves, measured
exactly as it would be on a flight. We extract every crest line of the French
Alps and of the French Pyrenees from a 25~m digital elevation model (DEM),
join the pieces into long paths that follow the main ridges, and let a ball
run along each path at a constant speed. The ball is a pilot with no pauses
and no choices, and its time-averaged mean squared displacement (TAMSD) gives
a geometric exponent, $\Hgeo$, that depends on the crest shape only.

Three findings come out. First, $\Hgeo$ is close to one but not one: about 0.92
over lags from 100~m to tens of kilometres. Second, it is scale dependent,
falling steadily from 0.98 to below 0.9 with no plateau, so a single number is
a summary and the local exponent $H(s)$ is the real result. Third, it is almost
the same in the two massifs, which have very different structures: the Alps
are a wide band of parallel chains, the Pyrenees a single long divide.
Box counting confirms the TAMSD result independently.

The paper is the first half of a comparison. Its companion (on real flights,
with the code and data shipped here) will put the tracks of paragliders on the
same axes and split their exponent into terrain, route and time. We only
sketch that step in the discussion.

\section{Crest lines from a DEM}
\label{sec:crests}

\paragraph{Data and region.}
We use the IGN BD ALTI elevation model at 25~m resolution \citep{ign_bdalti},
in Lambert-93. A massif is the set of departments of its loi Montagne massif,
nine in the Alps and six in the Pyrenees (Appendix~\ref{app:params}); within
them we keep the mountain area of the UNEP-WCMC definition \citep{kapos2000},
evaluated on 1~km cells from the DEM itself. The mountain area is
47\,700~km$^2$ in the Alps and 20\,600~km$^2$ in the Pyrenees
(Figure~\ref{fig:map}; Appendix~\ref{app:pyrenees}). The DEM stops at the
national border, which in the Pyrenees and in much of the Alps runs along the
main divide; we flag the crests that reach the edge of the data and repeat
every result without them (Appendix~\ref{app:edge}).

\paragraph{Crests are divides.}
A crest is where the ground falls away on both sides: a drop of water just
on one side of it ends in one valley, a drop on the other side in another
(Figure~\ref{fig:schematic}a). We therefore find crests as the boundaries
between the areas drained by different valleys. The DEM is depression-filled
with a priority flood \citep{barnes2014}, every cell drains to its steepest
neighbour, and the valley network is the set of cells draining at least
0.5~km$^2$, cut into links at the confluences. Two neighbouring cells that
drain into different links sit on a crest, unless one link is downstream of
the other. Divides are thus continuous through summits and cols, as crest
lines should be. Two tests remove divides that are not crests: the ground
150~m away on each side must be at least 10~m lower, and the crest must stand
at least 20~m above the streams on both sides, which rejects divides that run
across valley floors.

\paragraph{From a network to curves.}
The crest cells' sides form a planar network. At each junction a curve goes on
along the straightest continuation (deflection at most $135^\circ$), as a pilot
who keeps following the ridge would, and the other branches become curves of
their own. Flat stretches up to 1~km are bridged, dead ends shorter than 100~m
are pruned, and curves shorter than 1~km are dropped. Each curve is then
smoothed by a spline and resampled every $\Delta s = 5$~m of arc length
(Figure~\ref{fig:schematic}b). The Alps give 15\,205 curves (42\,000~km in
total, median 1.8~km), the Pyrenees 6\,707 (19\,500~km, median 1.9~km).

\paragraph{Joined paths.}
These curves are short: a ridge breaks wherever it flattens or where the
divide jumps to a parallel line. A pilot does not stop there, so we join the
pieces (Figure~\ref{fig:schematic}c). From every dead end we look for a
straight join of at most 1~km to another piece of the crest network, within
$45^\circ$ of the direction in which the curve runs out, that never drops
below 20~m above the streams: the ball does not cross a valley. Joins are made
shortest first, each dead end taking part in one join at most, so the pieces
grow into trees. A tree's root is one end of its longest path, and its paths
run from the root to every other end. The Alps give 1\,413 trees and 16\,096
paths (median 40~km, longest 309~km), the Pyrenees 483 trees and 7\,154 paths
(median 62~km, longest 252~km). Paths of one tree share their trunk, so every
bootstrap below resamples whole trees.

\begin{figure}[t]
\centering
\includegraphics[width=\textwidth]{schematic}
\caption{Method, stylized (not data).
\textbf{a}~A crest is the divide between two valleys: water released on either
side of it drains into a different stream.
\textbf{b}~Divides are found on the sides of the 25~m DEM cells (grey steps) and
smoothed by a spline (black). Below about four cells, 100~m, the shape of a
curve is set by the grid and the spline, not by the terrain.
\textbf{c}~Crest curves broken at flat cols are joined into trees (red: a join);
a path runs from the root to a tip (blue: one of the four paths of this tree).
\textbf{d}~A ball runs along a path at constant speed $v$; the TAMSD
$\delta^2(s)$ is the mean squared distance between its positions an arc length
$s = v\tau$ apart.}
\label{fig:schematic}
\end{figure}

\begin{figure}[t]
\centering
\includegraphics[width=\textwidth]{map_alps}
\caption{The crest curves of the French Alps, from the 25~m IGN DEM.
\textbf{a}~The whole massif inside its mountain area (black outline), with
the curves shaded by importance, the area of the smaller valley a curve
separates, and the gap joins in red.
\textbf{b}~Detail of the \'Ecrins.
\textbf{c}~Detail of the Haute-Maurienne, where the data stop at the border
and the curves reaching the edge (orange) are flagged
(Appendix~\ref{app:edge}). The Pyrenees are in Appendix~\ref{app:pyrenees}.}
\label{fig:map}
\end{figure}

\section{A ball on the crest}
\label{sec:tamsd}

A ball runs along each curve or path at constant speed $v$. Since the curve is
sampled every $\Delta s$ of arc length, the ball is at $\mathbf r_i =
\mathbf r(i\,\Delta s)$ at time $t_i = i\,\Delta s/v$, and a time lag $\tau$ is
an arc length $s = v\tau$. For a curve of $N$ points the TAMSD is
\begin{equation}
  \delta^2(s = k\,\Delta s) = \frac{1}{N-k}\sum_{i=0}^{N-k-1}
  \left|\mathbf r_{i+k} - \mathbf r_i\right|^2 ,
  \label{eq:tamsd}
\end{equation}
the standard time average of single-particle tracking \citep{metzler2014}.
It does not depend on the origin nor on the direction of travel. The speed
only rescales the time axis; we use $v = 35$~km/h, a typical cross-country
speed, so 100~m is 10.3~s and 10~km is about 17 minutes, but every result is
really a function of $s$. A straight line has $\delta^2 = s^2$; any bending
makes $\delta^2 < s^2$.

The lags $k$ are 80 log-spaced values, and on each curve they go up to a
quarter of its length: longer lags rest on too few independent windows. At
each lag we average $\delta^2$ over curves with equal weight per curve, and use
only lags with at least 10 curves. The local exponent of the mean is
\begin{equation}
  H(s) = \frac12\,\frac{d\ln\overline{\delta^2}}{d\ln s},
  \label{eq:hlocal}
\end{equation}
and $\Hgeo$ is the least-squares slope of $\ln\overline{\delta^2}$ against
$\ln s$, divided by two, over a window of lags. Since the lags are
log-spaced, every decade weighs the same. The 95\% intervals come from 1000
bootstrap resamples of the curves, or of whole trees for the paths.

\paragraph{Where the fit starts.}
Below a few grid cells the shape of a crest is not the terrain's: the divide
runs on cell sides, and the spline that smooths these steps makes the curve
straight below its own scale (Figure~\ref{fig:schematic}b). The fit therefore
starts at four cells, $s = 100$~m, and every window below begins there or
later. Below 100~m the curves are almost straight by construction
($H(s) > 0.98$), and including those lags would push $H$ towards one.

\section{Results: the terrain alone}
\label{sec:results}

\begin{figure}[t]
\centering
\includegraphics[width=\textwidth]{geometric}
\caption{The geometric TAMSD of the crests of the French Alps (blue) and
Pyrenees (orange), for the joined paths (solid) and the single curves
(dashed). Bottom axis: lag $\tau$ of a ball at 35~km/h; top axis: the same lag
as arc length. Grey: below four DEM cells, where the shape is set by the grid.
\textbf{a}~Mean TAMSD over its value on a straight line.
\textbf{b}~Local Hurst exponent, Eq.~\eqref{eq:hlocal}, with its 95\% interval
(bootstrap over trees; shown for the paths).
\textbf{c}~$\Hgeo$ fitted on windows of arc length, with 95\% intervals (filled:
paths; open: single curves, which are too short for the last window).}
\label{fig:geometric}
\end{figure}

Figure~\ref{fig:geometric} and Table~\ref{tab:hgeo} hold the main result.
Over the whole range from 100~m, the joined paths have $\Hgeo = 0.922$ in the
Alps and $0.916$ in the Pyrenees; the single curves, which are shorter and
straighter, $0.944$ and $0.939$. The difference between the massifs is within
the bootstrap intervals for the paths and about 0.005 for the curves.

The single number hides the main feature: the exponent is scale dependent.
The local exponent of the Alps paths is 0.987 at 100~m, 0.973 at 230~m, 0.949
at 550~m, 0.94 at 2~km, 0.89 at 4.6~km and 0.87 at 11~km
(Figure~\ref{fig:geometric}b). There is no plateau. Fitted on windows, $\Hgeo$
goes from 0.96 between 0.1 and 1~km to 0.925 between 1 and 10~km and about
0.86--0.87 beyond 10~km (Figure~\ref{fig:geometric}c). The same happens in
the Pyrenees, along nearly the same curve.

The ratio $\overline{\delta^2}/s^2$ (Figure~\ref{fig:geometric}a) gives the same
information in a more direct form. Its square root is the end-to-end distance
over the arc length: about 0.9 at 1~km and about 0.6 at 50~km. A ball that
follows the ridges for 50~km of arc gets about 30~km away from where it
started.

\begin{table}[t]
\centering
\caption{$\Hgeo$ fitted on windows of arc length $s$, with 95\% bootstrap
intervals (over curves, or over whole trees for the paths). The single curves
are too short for the last window.}
\label{tab:hgeo}
\small
\setlength{\tabcolsep}{4pt}
\begin{tabular}{lllll}
\toprule
 & \multicolumn{2}{c}{Alps} & \multicolumn{2}{c}{Pyrenees} \\
\cmidrule(lr){2-3}\cmidrule(lr){4-5}
window of $s$ & paths & curves & paths & curves \\
\midrule
from 0.1~km & 0.922 (0.918--0.929) & 0.944 (0.940--0.948) & 0.916 (0.910--0.931) & 0.939 (0.933--0.943) \\
0.1--1~km   & 0.960 (0.959--0.961) & 0.965 (0.964--0.965) & 0.957 (0.953--0.960) & 0.964 (0.963--0.965) \\
1--10~km    & 0.925 (0.920--0.933) & 0.928 (0.919--0.936) & 0.910 (0.900--0.925) & 0.912 (0.894--0.925) \\
beyond 10~km & 0.868 (0.773--0.877) & -- & 0.859 (0.841--0.946) & -- \\
\bottomrule
\end{tabular}
\end{table}

The intervals beyond 10~km are wide and asymmetric. The paths are many, but
the long lags rest on a few big trees. A lag of 10~km needs a path of at
least 40~km, since lags stop at a quarter of a path, and only 32 of the
1\,413 trees of the Alps have one, 12 of the 483 in the Pyrenees. A
bootstrap that draws whole trees sees whether one of them is in or out, so
the last window carries more uncertainty than its number of paths suggests.

\section{Box counting as an independent check}
\label{sec:box}

The TAMSD measures how far the ball gets; box counting measures how much room
the curve takes. For a curve parametrized by arc length with $R(s) =
\sqrt{\delta^2(s)} \sim s^H$, boxes of size $\varepsilon = R(s)$ each hold a
stretch of length about $s$, so a curve of length $L$ needs $N(\varepsilon)
\approx L/s \sim \varepsilon^{-1/H}$ of them. The box dimension of the curve is
then
\begin{equation}
  D = 1/H,
  \label{eq:dh}
\end{equation}
the relation of self-similar curves \citep{mandelbrot1967,mandelbrot1985}, as
opposed to $D = 2 - H$ for the graph of a self-affine profile. More generally,
$D(\varepsilon)$ at $\varepsilon = R(s)$ should equal $1/H(s)$.

The difficulty is finite size. A box count over a population of curves of
different lengths saturates at both ends: below the sampling step every point
is its own box, and beyond the length of a curve one box holds all of it. We
take this out with a control: each curve or path is matched with a straight
segment of the same arc length, first point and end-to-end direction, counted
on the same grids. The ratio of the mean counts $N/N_{\mathrm{straight}}$ keeps
only what the shape of the crests changes, and
\begin{equation}
  D(\varepsilon) = 1 + \frac{d \ln (N/N_{\mathrm{straight}})}{d\ln(1/\varepsilon)}.
\end{equation}
Every curve is counted on four rotated grids at 73 box sizes from 1~m to
1000~km, and the 95\% interval of $D$ comes from a bootstrap over curves, or
over trees for the paths.

\begin{figure}[t]
\centering
\includegraphics[width=\textwidth]{boxcount}
\caption{Box counting. \textbf{a}~Mean number of boxes against the inverse box
size, for the crests and for straight segments of the same lengths (paths
only; the curves' controls overlap the curves). Grey: boxes smaller than four
DEM cells. \textbf{b}~Ratio of the counts of crests and of straight segments.
\textbf{c}~Local Hurst exponent of the TAMSD of the paths (solid, as in
Figure~\ref{fig:geometric}b) and $1/D$ at $\varepsilon = R(s)$ (dotted), with
95\% intervals.}
\label{fig:box}
\end{figure}

Figure~\ref{fig:box} shows the result. The two saturations are clear in
panel~a, and the ratio in panel~b removes them: for the paths it dips to
0.7--0.75 at boxes of tens of kilometres, for the single curves only to 0.95
at about 2~km, since they are short and nearly straight. The box count of the
single curves carries information only up to about 1~km.

For the joined paths, $1/D$ follows $H(s)$ (panel~c). In the Pyrenees it is
0.990 against 0.982 at 100~m, 0.959 against 0.946 at 750~m and 0.935 against
0.926 at 3.4~km; in the Alps it is 0.01--0.02 above $H(s)$ over the same
range. Beyond about 10~km the box count of both massifs drifts up towards one
while $H(s)$ keeps falling: a few long paths fill the largest boxes, and the
finite-size control can only take out a straight line's share of that. Up to
10~km the two measures, built on different averages, agree on the shape of
the crests to within 0.01--0.03, with $1/D$ slightly higher.

\section{Discussion and outlook}
\label{sec:discussion}

\paragraph{The baseline.}
A ball at constant speed along the crests of the Alps or the Pyrenees has a
Hurst exponent of about 0.96 between 0.1 and 1~km, 0.92 between 1 and 10~km,
and below 0.9 beyond. These are not small corrections. A flight track with
$H = 0.92$ on the 1--10~km window is exactly as straight as the ridges, and
departures from ballistic motion of that size need no explanation in terms of
the pilot. That the two massifs, so different in layout, give the same curve
suggests that the result is a property of eroded mountain ridges in general,
in the spirit of the scale invariance of river networks
\citep{rodriguez1997}; we have only two massifs, so this remains a hint.

\paragraph{Terrain, route, time.}
A real flight differs from the ball in two ways (Figure~\ref{fig:outlook}).
The pilot's route is not the crest: it takes shortcuts across valleys, flies
beside the ridge, and circles in thermals. And the motion along the route is
not uniform: glides alternate with long climbs in which the pilot circles and
makes little progress. This suggests three exponents measured with the same
TAMSD and on the same windows of arc length:
\begin{itemize}
\item $\Hgeo(s)$, the crests, this paper;
\item $\Hroute(s)$, the pilot's track resampled by arc length with the
  thermalling circles replaced by their chords;
\item $\Hdyn(\tau)$, the track in time, as it is recorded.
\end{itemize}
$\Hroute - \Hgeo$ measures how the pilot's choices bend or straighten the
path relative to the ridges; $\Hdyn$ against $\Hroute$ measures what the
pauses do. If the progress along the route over a lag $\tau$ grows as
$\tau^\beta$, the displacement in time is that of the route evaluated at that
progress, so
\begin{equation}
  \Hdyn(\tau) \approx \beta(\tau)\,\Hroute\bigl(s(\tau)\bigr),
\end{equation}
a subordination in the sense of walks with waiting times
\citep{havlin1987,metzler2014}. When the pauses have a finite mean, as
expected for thermal climbs, $\beta$ is one at lags much longer than a glide
and a climb, and the pauses act only on the crossover between the two;
heavy-tailed pauses would lower $\beta$ over a wide range of lags. Which of
the two holds is a question for the data.

\begin{figure}[t]
\centering
\includegraphics[width=\textwidth]{outlook}
\caption{The decomposition proposed for real flights, stylized.
\textbf{a}~The crest sets $\Hgeo(s)$.
\textbf{b}~The pilot's route follows it with shortcuts and thermalling circles;
with the circles replaced by their chords it gives $\Hroute(s)$.
\textbf{c}~In time the pilot alternates progress (gliding) with pauses
(circling), which gives $\Hdyn(\tau)$.}
\label{fig:outlook}
\end{figure}

\paragraph{A test of the flight analysis.}
The folder that comes with this paper contains a reference implementation of
the three measurements for IGC flight files (\code{code/flights.py}) and a
test on synthetic flights. A synthetic pilot flies 60 of the Alps paths at
35~km/h, with glides and climbs of random durations, circles of 40~m radius in
the climbs and, in one variant, a 3~m correlated GPS error and two minutes on
the ground at each end. The route exponent comes back within 0.004 of the
paths' own $\Hgeo$ (0.937--0.940 against 0.939 from 100~m), with or without
noise. Without the circles collapsed, the track has $H = 0.906$: the circles
alone lower $H$ by 0.03, mostly below 1~km. Thermalling circles must
therefore be removed before a flight is compared with the terrain.

\paragraph{Limits.}
The crest network is a proxy for the ridges pilots follow, not a map of them.
Its shape below 100~m is the grid's, which is why the fits start there. The
joins across gaps, though restricted to 1~km and never through a valley,
decide which pieces form a path, and the longest lags rest on a few big trees.
The DEM stops at the border, and in both massifs the main divide is partly
outside it; dropping the curves that reach the edge of the data leaves every
exponent unchanged to the third decimal (Appendix~\ref{app:edge}), but the
longest paths are still cut there. Finally, 35~km/h only labels the time axis:
every comparison with flights should be made in arc length.

\section*{Data and code}
Everything needed to reproduce the figures from the DEM, and to run the
flight analysis, is in the folder of this paper: \code{code/} (the pipeline and
\code{flights.py}), \code{data/} (the per-curve TAMSD, the window fits, the box
counts, the crests, the masks) and \code{FLIGHTS.md}, a recipe for the
comparison with real flights.

\begin{thebibliography}{9}

\bibitem[Barnes et al.(2014)]{barnes2014}
Barnes, R., Lehman, C. and Mulla, D. (2014).
Priority-flood: an optimal depression-filling and watershed-labeling
algorithm for digital elevation models.
\emph{Computers \& Geosciences} 62, 117--127.

\bibitem[Havlin and Ben-Avraham(1987)]{havlin1987}
Havlin, S. and Ben-Avraham, D. (1987).
Diffusion in disordered media.
\emph{Advances in Physics} 36, 695--798.

\bibitem[IGN(2021)]{ign_bdalti}
IGN (2021).
\emph{BD ALTI version 2.0, 25~m}. Institut national de l'information
g\'eographique et foresti\`ere, Licence Ouverte 2.0.

\bibitem[Kapos et al.(2000)]{kapos2000}
Kapos, V., Rhind, J., Edwards, M., Price, M. F. and Ravilious, C. (2000).
Developing a map of the world's mountain forests.
In Price, M. F. and Butt, N. (eds.), \emph{Forests in Sustainable Mountain
Development}, CABI, Wallingford, 4--9.

\bibitem[Mandelbrot(1967)]{mandelbrot1967}
Mandelbrot, B. (1967).
How long is the coast of Britain? Statistical self-similarity and fractional
dimension.
\emph{Science} 156, 636--638.

\bibitem[Mandelbrot(1985)]{mandelbrot1985}
Mandelbrot, B. B. (1985).
Self-affine fractals and fractal dimension.
\emph{Physica Scripta} 32, 257--260.

\bibitem[Metzler et al.(2014)]{metzler2014}
Metzler, R., Jeon, J.-H., Cherstvy, A. G. and Barkai, E. (2014).
Anomalous diffusion models and their properties: non-stationarity,
non-ergodicity, and ageing at the centenary of single particle tracking.
\emph{Physical Chemistry Chemical Physics} 16, 24128--24164.

\bibitem[Rodr\'iguez-Iturbe and Rinaldo(1997)]{rodriguez1997}
Rodr\'iguez-Iturbe, I. and Rinaldo, A. (1997).
\emph{Fractal River Basins: Chance and Self-Organization}.
Cambridge University Press.

\end{thebibliography}

\appendix

\section{Extraction parameters}
\label{app:params}

\begin{table}[h]
\centering
\caption{Parameters of the crest extraction (\code{code/build.py},
\code{code/joined.py}, \code{code/mountains.py}).}
\begin{tabularx}{\textwidth}{lX}
\toprule
DEM & IGN BD ALTI v2, 25~m, Lambert-93 \\
Alps & departments 04, 05, 06, 26, 38, 73, 74, 83, 84 \\
Pyrenees & departments 09, 11, 31, 64, 65, 66 \\
Mountain area & UNEP-WCMC classes \citep{kapos2000} on 1~km cells, local
  elevation range within 7~km, enclosed basins up to 25~km$^2$ added, largest
  connected piece kept; a curve belongs to the massif if at least half of it
  lies in this area \\
Valleys & cells draining at least 0.5~km$^2$ \\
Cross-section & ground 6 cells (150~m) away on each side at least 10~m lower \\
Height above streams & at least 20~m on both sides \\
Gaps bridged & flat stretches up to 40 cell sides (1~km) \\
Pruning & dead ends shorter than 4 cell sides (100~m) \\
Junctions & a stroke goes on if it deflects by at most $135^\circ$, direction
  measured 200~m from the junction \\
Curves & spline, resampled every 5~m; shorter than 1~km dropped \\
Joins & straight, at most 1~km, within $45^\circ$ of the direction the curve
  runs out, never lower than 20~m above the streams, snapped to the other
  curve's dead end within 200~m, up to 3 components tried per dead end \\
Edge of the data & within 4 cells (100~m) of a cell without data \\
TAMSD & 80 log-spaced lags up to a quarter of each curve, lags with at least
  10 curves, fit from 100~m, 1000 bootstrap resamples \\
Box counting & 73 sizes from 1~m to 1000~km, 4 rotated grids, 400 bootstrap
  resamples, local slope over half a decade \\
\bottomrule
\end{tabularx}
\end{table}

The loi Montagne is used only as a list of departments: the mountain area
itself is recomputed from the DEM.

\section{The Pyrenees}
\label{app:pyrenees}

\begin{figure}[h]
\centering
\includegraphics[width=\textwidth]{map_pyrenees}
\caption{The crest curves of the French Pyrenees, as in Figure~\ref{fig:map}.}
\label{fig:map_pyrenees}
\end{figure}

\section{Curves at the edge of the data}
\label{app:edge}

Without the curves and paths that reach the edge of the data (within 100~m of
a cell without data), the exponents from 100~m are: Alps paths 0.922
(0.918--0.931), Alps curves 0.944 (0.940--0.948), Pyrenees paths 0.916
(0.910--0.931), Pyrenees curves 0.939 (0.933--0.943). They are the values of
Table~\ref{tab:hgeo} to the third decimal.

\end{document}
